% AUTORIDAD DE FUENTE:
% expose_VIII_body.tex, líneas 1701--1743; índices de escaneo 284--286;
% folios de fuente 273--275. La referencia impresa (1.11.1) y su
% corrección local (4.11.1) se conservan reversiblemente mediante \sisi.

%%%%%%%%%%  scan idx 284  |  folio 273  %%%%%%%%%%

\medskip
\noindent 4.11. Sean $P$ y $Q$ grupos algebraicos lisos sobre $k$,
sea $E$ una biextensión de $(P,Q)$ por $\underline G_m$, y sea $T$ el
mayor subtoro de $P$. Determinaremos el anulador izquierdo del
emparejamiento
\begin{equation}\tag{4.11.1}
 T_\ell(P)\times T_\ell(Q)
 \longrightarrow T_\ell(\underline G_{m,k}),
\end{equation}
donde $\ell$ vuelve a ser un número primo fijo. Por 4.10, este anulador
contiene siempre el subgrupo $T_\ell(T)$ de $T_\ell(P)$.

Elijamos una extensión $k'$ de $k$ sobre la cual $P$ y $Q$ puedan
escribirse cada uno como una extensión de un esquema abeliano por un
grupo afín que sea una extensión iterada de toros y copias de
$\underline G_a$. Si $P$ y $Q$ tienen rango unipotente cero, podemos
tomar $k'=k$; en general, podemos tomar como $k'$ la clausura perfecta
de $k$ (o una extensión radicial finita). Sean $A'$ y $B'$ las partes
abelianas de $P_{k'}$ y $Q_{k'}$, respectivamente. Por 4.7, la
biextensión $E_{k'}$ de $(P_{k'},Q_{k'})$ procede de una biextensión
$F'$ de $(A',B')$, determinada de manera única salvo isomorfismo único.
En consecuencia, el emparejamiento
\begin{equation}\tag{4.11.2}
 T_\ell(P_{k'})\times T_\ell(Q_{k'})
 \longrightarrow T_\ell(\underline G_{m,k'}),
\end{equation}
obtenido de (4.11.1) mediante la extensión del cuerpo base $k\to k'$,
y definido asimismo por la biextensión $E_{k'}$, se obtiene del
emparejamiento
\begin{equation}\tag{4.11.3}
 T_\ell(A')\times T_\ell(B')
 \longrightarrow T_\ell(\underline G_{m,k'})
\end{equation}
definido por $F'$, componiéndolo con los homomorfismos canónicos (que
son epimorfismos si $P$ y $Q$ son $\ell$-divisibles)
\begin{equation}\tag{4.11.4}
 T_\ell(P_{k'})\longrightarrow T_\ell(A'),
 \qquad
 T_\ell(Q_{k'})\longrightarrow T_\ell(B').
\end{equation}

%%%%%%%%%%  scan idx 285  |  folio 274  %%%%%%%%%%

Se sigue que el anulador izquierdo de (4.11.2), obtenido del anulador
izquierdo de (4.11.1) por cambio de base, es la imagen inversa, por el
primer homomorfismo de (4.11.4), del anulador izquierdo de (4.11.3).
Sea $C'$ el mayor subesquema abeliano de $A'$ tal que la restricción de
$F'$ a $(C',B')$ sea la biextensión trivial. Equivalentemente, $C'$ es
la componente neutra reducida del núcleo del homomorfismo
\begin{equation}\tag{4.11.5}
 A'\longrightarrow(B')^*
\end{equation}
definido por $F'$, donde $(B')^*$ es el esquema abeliano dual de $B'$.
El anulador izquierdo de (4.11.3) se identifica con $T_\ell(C')$; por
tanto, el anulador izquierdo de
$(\sisi{1.11.1}{4.11.1})$ es precisamente la imagen inversa de
\[
 T_\ell(C')\lhook\joinrel\longrightarrow T_\ell(A')
\]
por el primer morfismo de (4.11.4).

En particular, decir que el anulador izquierdo de (4.11.1) es
exactamente $T_\ell(T)$ —equivalentemente, que el anulador izquierdo de
(4.11.2) es el núcleo $T_\ell(T')$ del primer morfismo de (4.11.4)—
equivale a decir que (4.11.3) es no degenerado por la izquierda, es
decir, que $C'=0$, o, equivalentemente, que (4.11.5) tiene núcleo
finito. En esta última forma, la condición es manifiestamente
independiente de $\ell$, y en cualquiera de sus formas es claramente
independiente de la elección de $k'$. Cuando se cumple, diremos
sencillamente que la biextensión dada $E$ de $(P,Q)$ por
$\underline G_m$ es \underline{no degenerada por la izquierda}. Por
dualidad, esto significa también que el homomorfismo
\begin{equation}\tag{4.11.6}
 B'\longrightarrow(A')^*
\end{equation}
definido por $F'$ es un epimorfismo, es decir, es sobreyectivo.

%%%%%%%%%%  scan idx 286  |  folio 275  %%%%%%%%%%

La noción de biextensión \underline{no degenerada por la derecha} se
define simétricamente: (4.11.6) debe tener núcleo finito o,
equivalentemente, (4.11.5) debe ser un epimorfismo. Se dice que la
biextensión $E$ de $(P,Q)$ por $\underline G_{m,k}$ es
\underline{no degenerada} si lo es por la izquierda y por la derecha.
Esto equivale a decir que (4.11.5), o equivalentemente (4.11.6), es una
isogenia.

De la definición se sigue inmediatamente que estas nociones son
invariantes por toda extensión del cuerpo base $k$.
